%! Minimizing Loss by Gradient Descent — Unit 8, Lesson 8.3 — Guided Notes (Answer Key)
\documentclass[10pt]{article}
\usepackage{linalg-article}
\usepackage{linalg-key}   % pulls in -boxes; do NOT also load -boxes
\newcommand{\TallMath}[1]{$\displaystyle #1\rule[-1.4em]{0pt}{3.2em}$}
\newcommand{\vv}{\mathbf{v}}
\newcommand{\bb}{\mathbf{b}}
\newcommand{\ww}{\mathbf{w}}
\newcommand{\zero}{\mathbf{0}}
\newcommand{\relu}{\mathrm{ReLU}}

\begin{document}

\pageheader{Unit 8, Lesson 8.3}{Guided Notes --- Answer Key}
\namedateperiod

\vspace{0.12in}

\begin{objectivebox}
\textbf{By the end of this lesson, I will be able to\dots}
\begin{itemize}\setlength{\itemsep}{2pt}
  \item read a \emph{loss} $L(w)$ as a bowl and see that learning means finding its lowest point (the smallest error);
  \item use the \emph{slope} to tell which way is downhill, and take a gradient-descent step $w\leftarrow w-s\,L'(w)$;
  \item extend the same rule to many weights with the \emph{gradient} $\nabla L$, and explain how the learning rate $s$ controls the descent.
\end{itemize}
\end{objectivebox}

\vspace{0.06in}

\begin{vocabbox}
\small
Fill in each term as we build it together.
\vspace{2pt}
\termblanklong{Loss $L$ --- total (squared) error; smaller means a better fit}
\termblanklong{Gradient descent --- repeatedly step downhill on the loss}
\termblanklong{Slope $L'(w)$ --- points uphill; step the opposite way}
\termblanklong{Learning rate $s$ --- the size of each downhill step}
\termblanklong{Gradient $\nabla L$ --- the slopes for all weights at once}
\end{vocabbox}

\begin{hookbox}
We can now \emph{build} networks --- layers $\relu(A\vv+\bb)$ (8.0), folds (8.1), filters (8.2). But every choice
of weights gives a different \textbf{loss}: the Unit-4 sum of squared errors between the network's predictions
and the targets. Picture that loss as a landscape over the weights.
\begin{itemize}\setlength{\itemsep}{1pt}
  \item If ``good weights'' means ``small loss,'' where in the landscape do we want to end up? \ans{At the \textbf{lowest point} --- the bottom of the valley, where the loss is smallest.}
  \item You are in a valley in thick fog and can feel only the \emph{slope} under your feet. How would you walk to the bottom? \ans{Feel which way is downhill and take a step that way; repeat until the ground is flat.}
\end{itemize}
\end{hookbox}

\begin{notesbox}{1. Learning Means Minimizing the Loss}
The \textbf{loss} $L$ scores a set of weights: it is the Unit-4 \emph{sum of squared errors} between what the
network predicts and the targets. Big loss $=$ bad fit; small loss $=$ good fit. So \textbf{learning} is one
job:
\begin{center}
\emph{choose the weights that make the loss $L$ as \textbf{small} as possible.}
\end{center}
For a single weight $w$, the loss is a \textbf{bowl} (a parabola). Our example bowl is
\[
  L(w)=(w-4)^2,
\]
which is $0$ when $w=4$ (a perfect fit) and grows as $w$ moves away. The bottom of the bowl is the best
weight. The trouble: a real network has \emph{millions} of weights, so we cannot just eyeball the bottom --- we
have to \emph{walk} there.
\end{notesbox}

\begin{notesbox}{2. Which Way Is Downhill? The Slope}
At any weight $w$, the \textbf{slope} $L'(w)$ measures how steeply the loss is rising --- it points
\emph{uphill}. To make the loss \emph{fall}, step the \emph{opposite} way. (We give you the slope; no calculus
needed.) For our bowl the slope is
\[
  L'(w)=2w-8.
\]
At $w=0$: $\;L'(0)=2(0)-8={\color{keyred}\mathbf{-8}}$. The slope is \emph{negative}, so the loss is still
falling as $w$ \emph{increases} --- downhill is to the \textbf{right} (toward $w=4$). At the very bottom $w=4$,
the slope is $L'(4)={\color{keyred}\mathbf{0}}$: flat ground, nowhere lower to go.

\begin{center}
\begin{tikzpicture}[scale=0.92]
  % LEFT: the bowl + ball + downhill arrow (the idea)
  \begin{scope}[shift={(0,0)}]
    \draw[->, charcoal, thin] (-0.3,0)--(5.2,0) node[right]{\scriptsize $w$};
    \draw[->, charcoal, thin] (0,-0.25)--(0,2.1) node[above]{\scriptsize $L(w)$};
    \draw[ultra thick, royalblue, smooth]
      plot coordinates {(0,1.44)(0.6,0.81)(1.2,0.36)(1.8,0.09)(2.4,0)(3.0,0.09)(3.6,0.36)(4.2,0.81)(4.8,1.44)};
    \fill[burgundy] (0.6,0.81) circle (3pt);
    \draw[charcoal, thick] (0.15,1.24)--(1.05,0.38);
    \node[charcoal] at (0.35,1.55) {\scriptsize slope (uphill)};
    \draw[->, very thick, goldacc] (1.05,0.7)--(2.15,0.12);
    \node[goldacc] at (2.35,0.55) {\scriptsize step downhill};
    \fill[charcoal] (2.4,0) circle (1.6pt);
    \node[charcoal] at (3.05,-0.32) {\scriptsize bottom $=$ best $w$};
  \end{scope}
  % RIGHT: the descent staircase (the algorithm)
  \begin{scope}[shift={(6.4,0)}]
    \draw[->, charcoal, thin] (-0.3,0)--(5.2,0) node[right]{\scriptsize $w$};
    \draw[->, charcoal, thin] (0,-0.25)--(0,2.1) node[above]{\scriptsize $L(w)$};
    \draw[ultra thick, royalblue, smooth]
      plot coordinates {(0,1.44)(0.6,0.81)(1.2,0.36)(1.8,0.09)(2.4,0)(3.0,0.09)(3.6,0.36)(4.2,0.81)(4.8,1.44)};
    \fill[burgundy] (0,1.44) circle (2.6pt);   \node[burgundy] at (-0.25,1.62) {\scriptsize $w_0$};
    \fill[burgundy] (1.2,0.36) circle (2.6pt); \node[burgundy] at (1.2,0.66) {\scriptsize $w_1$};
    \fill[burgundy] (1.8,0.09) circle (2.6pt); \node[burgundy] at (1.86,0.36) {\scriptsize $w_2$};
    \fill[burgundy] (2.1,0.02) circle (2.6pt);
    \draw[->, goldacc, thick] (0.2,1.28)--(1.05,0.5);
    \draw[->, goldacc, thick] (1.32,0.28)--(1.72,0.14);
    \node[charcoal] at (3.1,0.95) {\scriptsize each step lands lower};
    \node[charcoal] at (3.05,-0.32) {\scriptsize $\to$ the minimum};
  \end{scope}
\end{tikzpicture}
\end{center}
\end{notesbox}

\begin{notesbox}{3. Gradient Descent: Step Against the Slope}
Put the idea into a rule. Start somewhere, and repeatedly take a step \emph{against} the slope:
\[
  \boxed{\,w \;\leftarrow\; w - s\,L'(w)\,}
\]
The number $s>0$ is the \textbf{learning rate} (step size). Roll down our bowl from $w=0$ with $s=0.25$:
\begin{center}
\renewcommand{\arraystretch}{1.35}
\begin{tabular}{c|c|c|c}
current $w$ & slope $L'(w)=2w-8$ & step $w-s\,L'(w)$ & new $w$\\ \hline
$0$ & $-8$ & $0-0.25(-8)$ & ${\color{keyred}\mathbf{2}}$\\
$2$ & $-4$ & $2-0.25(-4)$ & ${\color{keyred}\mathbf{3}}$\\
$3$ & $-2$ & $3-0.25(-2)$ & ${\color{keyred}\mathbf{3.5}}$\\
\end{tabular}
\end{center}
The weights march $0\to2\to3\to3.5\to\cdots$ straight toward the bottom at $w=4$, where the slope is $0$ and the
step changes nothing --- that is how we know we have arrived. The loss falls with every step: $16\to4\to1\to0.25$.
\end{notesbox}

\begin{notesbox}{4. Many Weights: The Gradient --- and Why It Matters}
A real network has millions of weights, so the loss is a bowl in millions of directions at once. The
\textbf{gradient} $\nabla L$ is just the list of slopes --- one for every weight --- packed into a vector. It
points in the steepest-\emph{uphill} direction, so we step \emph{against} it. The rule is \emph{exactly} the
same, now for the whole weight vector $\ww$:
\[
  \ww \;\leftarrow\; \ww - s\,\nabla L.
\]
For example, with $L(w_1,w_2)=(w_1-1)^2+(w_2-3)^2$ the gradient is $\nabla L=[\,2(w_1-1),\,2(w_2-3)\,]$. At
$(0,0)$ it is $[-2,-6]$, and one step with $s=0.25$ gives
\[
  (0,0)-0.25\,[-2,-6]=\big(\ {\color{keyred}\mathbf{0.5}},\ {\color{keyred}\mathbf{1.5}}\ \big),
\]
heading toward the minimum $(1,3)$.

\vspace{2pt}
\textbf{The learning rate matters.} Too small and the descent crawls; too big and a step overshoots the bottom
and can bounce \emph{uphill}. Choosing $s$ well is part of training.

\vspace{2pt}
\textbf{The whole picture.} Learning a network is Unit~1 (matrix layers) $+$ one bend (ReLU) $+$ Unit~4
(squared-error loss) $+$ rolling downhill (gradient descent) --- no new machinery, repeated millions of times.
\textbf{The road ahead:} \textbf{8.4} mean, variance, and covariance --- the shape of the data itself.
\end{notesbox}

\begin{practicebox}
Show your work. Use $L(w)=(w-4)^2$ with slope $L'(w)=2w-8$ and learning rate $s=0.25$ unless told otherwise.
\begin{enumerate}[leftmargin=1.6em, itemsep=4pt]
  \item Take one step from $w=6$. (Find $L'(6)$, then $w-s\,L'(6)$.) Which way did $w$ move, and toward what? \ansline{$L'(6)=2(6)-8=4$, so $w=6-0.25(4)=5$. It moved \emph{left}, toward the minimum $w=4$ --- from the right side this time.}
  \item Redo one step from $w=0$ but with a \emph{big} learning rate $s=1$. Where does $w$ land, and what went wrong? \ansline{$w=0-1(-8)=8$. The step \emph{overshot} the bottom ($w=4$) and jumped to the far wall; $L(8)=16=L(0)$, so no progress --- the learning rate was too big.}
  \item Two weights: $\nabla L=[\,2(w_1-1),\,2(w_2-3)\,]$. Take one step ($s=0.25$) from $(2,0)$. \ansline{$\nabla L=[\,2(2-1),\,2(0-3)\,]=[2,-6]$, so $(2,0)-0.25[2,-6]=(1.5,\,1.5)$ --- heading toward $(1,3)$.}
\end{enumerate}
\end{practicebox}

\begin{teachernote}
The one bowl $L(w)=(w-4)^2$ runs through the whole lesson (minimum $w=4$, slope $2w-8$). Keep three ideas
front and center: (1) the slope points \emph{uphill}, so we \emph{subtract} it; (2) the learning rate $s$ sets
the step size; (3) the slope is $0$ at the bottom, which is the stopping signal. The $\S3$ table is the core
skill --- $0\to2\to3\to3.5$ toward $w=4$, loss $16\to4\to1\to0.25$. Practice~2 is the ``too-big learning rate''
cautionary case (overshoot to $w=8$, same loss). $\S4$ and Practice~3 lift the identical rule to a gradient
vector; students often forget the minus sign or step \emph{with} the slope --- reinforce ``against the slope,''
i.e.\ downhill. Common slips: sign of the slope, and multiplying $s\,L'(w)$ before subtracting.
\end{teachernote}

\end{document}
